% !TeX root = ../../Book3.tex
% 纲要标题：### 10.4 数域整数环、理想类群与曲线赋值
% 纲要位置：工作记录/计划纲要.md，第 2126 行。
% 纲要细目（写作范围）：
% * 数域整数环的 Dedekind 性
% * 理想类群及 Picard 群的初步关系
% * 正规曲线与离散赋值
% ---


\section{数域整数环、理想类群与曲线赋值}
\label{b3:sec:1004}

整数环和一元多项式环都具有唯一的元素分解，但其有限扩张未必仍有这个性质。Dedekind 理论说明，转向非零理想之后仍能保持唯一分解。曲线上的零点阶数则把同一结构解释为几何数据。


\subsection{数域整数环的 Dedekind 性}
\label{b3:subsec:1004:01}

\begin{definition}\label{b3:def:number-field-integers}
给定有限域扩张 \(K/\mathbb Q\)，称 \(K\) 为\term[shuyu]{数域}[number field]。在 \(K\) 中，把所有满足某个首一整数系数多项式的元素组成的环称为 \(K\) 的\term[zhengshuhuan]{整数环}[ring of integers]，记为 \(\mathcal O_K\)；它也就是 \(\mathbb Z\) 在 \(K\) 中的整闭包。
\end{definition}
\symidx{\(\mathcal O_K\)：数域的整数环}

\begin{theorem}\label{b3:thm:number-integers-dedekind}
设 \(K/\mathbb Q\) 为有限域扩张，\(\mathcal O_K\) 为 \(\mathbb Z\) 在 \(K\) 中的整闭包。则 \(\mathcal O_K\) 是有限自由 \(\mathbb Z\) 模，秩为 \([K:\mathbb Q]\)，并且是 Dedekind 整环。
\end{theorem}
\begin{proof}
特征零保证 \(K/\mathbb Q\) 可分。\(\mathbb Z\) 是整闭 Noether 整环，故\cref{b3:lem:normalization-separable-finite}给出 \(\mathcal O_K\) 的有限生成性。
为直接说明自由性，令 \(n=[K:\mathbb Q]\)，选取有理基 \(\beta_1,\ldots,\beta_n\) 及坐标同构
\[
\theta:K\longrightarrow\mathbb Q^n,\qquad
\sum_{i=1}^n q_i\beta_i\longmapsto(q_1,\ldots,q_n).
\]
为 \(\mathcal O_K\) 的有限个整数模生成元的坐标取正整数公共分母 \(D\)，便有
\(\theta(\mathcal O_K)\subseteq D^{-1}\mathbb Z^n\)。映射
\[
\mathcal O_K\longrightarrow H\coloneqq D\theta(\mathcal O_K)\subseteq\mathbb Z^n,
\qquad a\longmapsto D\theta(a)
\]
是 \(\mathbb Z\) 模同构。每个子群 \(H\subseteq\mathbb Z^n\) 都自由，可对 \(n\) 归纳：第一坐标像若为零，就在 \(\mathbb Z^{n-1}\) 中使用归纳；否则其像为 \(d\mathbb Z\)、\(d>0\)，取 \(h\in H\) 的第一坐标为 \(d\)。每个元素唯一写成 \(mh+h_0\)，其中 \(m\in\mathbb Z\)、\(h_0\) 第一坐标为零，故 \(H=\mathbb Zh\oplus H_0\)，再对 \(H_0\) 归纳。这证明 \(\mathcal O_K\) 为有限自由模。
任何 \(x\in K\) 乘一个非零整数后都对 \(\mathbb Z\) 整：对 \(x\) 的首一有理系数方程清除分母，正如\cref{b3:lem:normalization-separable-finite}证明中的缩放。因此
\(\mathbb Q\otimes_{\mathbb Z}\mathcal O_K=K\)，自由秩为 \([K:\mathbb Q]\)，分式域也为 \(K\)。

\(\mathcal O_K\) 作为 \(\mathbb Z\) 模有限生成，因而是 Noether 环；整闭包对整性传递而整闭。
由\cref{b3:thm:integral-dimension}，它与 \(\mathbb Z\) 同为一维。于是它是 Dedekind 整环。
\end{proof}

这个证明也说明正规化有限性在算术中的作用：它先给出 Noether 性，不能只凭“所有元素都整”就断言整数环 Noether。
对任意 \(0\ne a\in\mathcal O_K\)，乘以 \(a\) 是有限自由整数模上的单射，且在 \(\mathbb Q\) 上为同构，所以余核是有限交换群。
因此 \(\mathcal O_K/(a)\) 有限；非零素理想的剩余域也有限。

\subsection{理想类群与 Picard 群}
\label{b3:subsec:1004:02}

在 Dedekind 整环上，非零分式理想都可逆，并由\cref{b3:prop:invertible-local-principal}成为局部自由秩一的有限生成投射模。若只保留其模结构，就得到一种在各局部环上都像一条自由直线的对象。这样的模在任意交换环上仍可用张量积相乘；Picard 群把它们的同构类组织起来。

\begin{definition}\label{b3:def:ideal-class-picard}
设 \(R\) 为 Dedekind 整环，\(K=\operatorname{Frac}(R)\)。分式理想群对其主分式理想子群的商
\[
\operatorname{Cl}(R)\coloneqq
\{\text{非零分式理想}\}/\{xR:x\in K^*\}
\]
称为 \(R\) 的\term[lixiangleiqun]{理想类群}[ideal class group]。
另设 \(R\) 为任意交换含幺环。考虑有限生成投射 \(R\) 模 \(L\) 的同构类，要求对每个素理想 \(\mathfrak p\subseteq R\)，\(L_{\mathfrak p}\) 都是秩一自由 \(R_{\mathfrak p}\) 模。这些同构类以 \(R\) 上的张量积为乘法所成的群，称为 \(R\) 的\term[Picard-qun]{Picard 群}[Picard group]，记为 \(\operatorname{Pic}(R)\)。
\end{definition}
\symidx{\(\operatorname{Cl}(R)\)：Dedekind 整环的理想类群}
\symidx{\(\operatorname{Pic}(R)\)：秩一有限投射模的 Picard 群}

说明后一构造确实为群。两个有限投射模都是有限自由模的直和因子，张量后仍为有限投射模；局部化后两个秩一自由模的张量积仍秩一。
对偶 \(L^*=\operatorname{Hom}_R(L,R)\) 也为有限投射模：将分解 \(R^r=L\oplus L'\) 取对偶便得到相应的直和分解。有限投射模也有限展示，因为这项分解中 \(L'\) 是由 \(R^r\) 的有限基经投影得到的有限生成核。因此\cref{b3:thm:hom-localization}适用于 \(L\)，它使求值映射
\(L\otimes_RL^*\rightarrow R\) 在每个素理想处成为秩一自由模的求值同构。
由\cref{b3:thm:local-morphism-detection}它本身即为同构。因此 \([L^*]\) 是逆元，\([R]\) 是单位元。

\begin{proposition}\label{b3:prop:dedekind-picard-class}
设 \(R\) 为 Dedekind 整环，\(K=\operatorname{Frac}(R)\)。则有自然群同构
\(\operatorname{Cl}(R)\simeq\operatorname{Pic}(R)\)，由 \([I]\mapsto[I]\) 给出。
而 \(\operatorname{Cl}(R)=0\) 当且仅当 \(R\) 是主理想整环。
\end{proposition}
\begin{proof}
可逆理想由\cref{b3:prop:invertible-local-principal}为有限投射模，并且局部自由秩一。
乘法映射 \(I\otimes_RJ\rightarrow IJ\) 在每个极大理想 \(\mathfrak m\) 处，把两个非零主分式理想 \(I_{\mathfrak m}=aR_{\mathfrak m}\)、\(J_{\mathfrak m}=bR_{\mathfrak m}\) 的张量积的基 \(a\otimes b\) 送到 \((IJ)_{\mathfrak m}\) 的基 \(ab\)，故局部为同构，再由局部检测知它是同构。
主分式理想同构于 \(R\)，所以映射在类群上良定义且保持乘法。

若有 \(R\) 模同构 \(\varphi:I\xrightarrow{\sim}J\)，借助自然乘法同构 \(K\otimes_RI\to K\)、\(\lambda\otimes a\mapsto\lambda a\)，以及 \(J\) 的对应同构，将 \(1\otimes\varphi\) 识别为 \(K\) 的一个线性自同构 \(\varphi_K\)。它必为乘以某个 \(x\in K^*\)。对每个 \(a\in I\)，自然嵌入 \(I,J\subseteq K\) 与张量映射相容，因而在 \(K\) 中有
\[
\varphi(a)=\varphi_K(a)=xa.
\]
所以 \(J=\varphi(I)=xI\)，说明类群映射单射。
给定有限投射且局部秩一的 \(L\)，它是自由模的直和因子，因而无挠；所以
\(L\rightarrow K\otimes_RL\) 单射。后者一维，选择它与 \(K\) 的同构，将 \(L\) 识别为 \(K\) 内有限生成的非零子模。为有限个生成元清除分母后，它成为一个分式理想。因此映射满射。
最后，类群平凡恰好表示每个非零分式理想都主生成，特别是每个非零通常理想主生成；反向由清分母同样成立。
\end{proof}

取 \(K=\mathbb Q(\sqrt{-5})\)，记 \(s=\sqrt{-5}\)。此时 \(R=\mathcal O_K=\mathbb Z[s]\)。
为核实这一步，设 \(\alpha=a+bs\) 对 \(\mathbb Z\) 整，其中 \(a,b\in\mathbb Q\)。它的共轭 \(\bar\alpha=a-bs\) 满足同一个首一整数系数多项式，因此也整。迹 \(\alpha+\bar\alpha=2a\) 和范数 \(\alpha\bar\alpha=a^2+5b^2\) 作为整元素的和与积也对 \(\mathbb Z\) 整，又都属于 \(\mathbb Q\)，由 \(\mathbb Z\) 整闭可知它们是整数。于是
\(-20b^2=(2a)^2-4(a^2+5b^2)\) 为整数。将 \(b\) 写成既约分数，分母的平方整除 \(20\)，所以 \(2b\in\mathbb Z\)。
写 \(a=A/2,b=B/2\)，由 \(A^2+5B^2\equiv0\pmod4\) 得 \(A,B\) 都为偶数，故 \(a,b\in\mathbb Z\)。反向由 \(s\) 整显然成立。

\ProseParagraphBreak
在 \(R=\mathbb Z[s]\) 中，范数 \(N(a+bs)=a^2+5b^2\) 是非负整数并具有乘法性。单位的范数必为 \(1\)，由这个公式可知单位恰为 \(\pm1\)。环中没有范数为 \(2\) 或 \(3\) 的元素：若 \(b\ne0\)，范数至少为 \(5\)；若 \(b=0\)，范数就是整数的平方。而
\[
N(2)=4,\qquad N(3)=9,\qquad N(1+s)=N(1-s)=6.
\]
若这些元素之一能分成两个非单位，其范数就会分成两个大于 \(1\) 的整数。但 \(4,9,6\) 的每种非平凡乘积分解都含因子 \(2\) 或 \(3\)，故 \(2,3,1+s,1-s\) 都不可约。相伴元素有相同范数，而 \(1+s\ne\pm(1-s)\)，所以这四个元素两两不相伴。因此
\[
6=2\cdot3=(1+s)(1-s)
\]
是两种不能由重排与乘单位互相得到的不可约分解，\(R\) 不是唯一分解整环。

\ProseParagraphBreak
理想 \(\mathfrak p=(2,1+s)\) 的商环为 \(\mathbb F_2\)，所以是素理想。
有 \(\mathfrak p^2=(2)\)：平方的生成元均被 \(2\) 整除，而
\(2(1+s)-(1+s)^2=6\)，连同 \(4\in\mathfrak p^2\) 得 \(2\in\mathfrak p^2\)。
它不是主理想。若 \(\mathfrak p=(\alpha)\)，则由 \(\mathfrak p^2=(2)\) 得 \(\alpha^2=2u\)，其中 \(u\in R^*=\{\pm1\}\)。取范数得 \(N(\alpha)^2=4\)，所以 \(N(\alpha)=2\)，与上述范数的取值矛盾。
因此 \([\mathfrak p]\) 是类群中阶二的非平凡元素。这里已经看到元素分解与理想分解的区别，不需要先计算整个类群。

\subsection{正规曲线与离散赋值}
\label{b3:subsec:1004:03}

设 \(k\) 为域，\(A\) 为有限型整 \(k\) 代数。若 \(\dim A=1\)，本节便称 \(A\) 所对应的仿射对象为一条不可约仿射曲线。这里用 \(\operatorname{Spec}A\) 表示该仿射对象，以便包含非代数闭基域的闭点。由\cref{b3:thm:affine-dimension}，其有理函数域在 \(k\) 上的超越次数为一。
若 \(A\) 整闭，则称曲线正规。此时 \(A\) 为 Dedekind 整环，每个闭点
\(\mathfrak p\) 的局部环是 DVR。

\ProseParagraphBreak
对 \(0\ne f\in\operatorname{Frac}(A)\)，整数 \(v_{\mathfrak p}(f)\) 表示在该点的阶数：正值是零点阶数，负值是极点阶数，零值表示在该局部环中为单位。
主分式理想的分解说明只有有限多个非零阶，并给出
\[
fA=\prod_{\mathfrak p\ne0}\mathfrak p^{v_{\mathfrak p}(f)}.
\]
这个分解的指数由 \(\operatorname{Spec}A\) 中的闭点给出。

\ProseParagraphBreak
在 \(A=k[t]\) 上，闭点由首一不可约多项式 \(q(t)\) 表示，统一化元为 \(q\)，阶数即 \(q\) 在分子和分母中的重数之差。
特别地，函数 \(t\) 在闭点 \((t)\) 处的阶数为 \(1\)，在其余闭点 \((q)\)（\(q\ne t\)）处的阶数都为 \(0\)。在同一函数域中令 \(u=t^{-1}\)，倒数坐标图 \(k[u]\) 中点 \((u)\) 的局部环是 \(k[u]_{(u)}\)，统一化元为 \(u\)。于是
\[
v_\infty(t)=v_{(u)}(u^{-1})=-1.
\]
这一点是在 \(t\) 坐标图之外补入的无穷远点。仿射环 \(k[t]\) 的理想分解记录了 \(t\) 在原点的一阶零点；倒数坐标则显示它在无穷远的一阶极点。

\ProseParagraphBreak
在 \(k[t,t^{-1}]\) 上，\((t)\) 已因局部化而消失，\(t\) 是单位，在留下的所有闭点的阶数都为零。
这解释了删去一点会怎样改变可逆函数，而不改变函数域。

\ProseParagraphBreak
尖点 \(k[t^2,t^3]\) 的正规化为 \(k[t]\)，见\cref{b3:prop:cusp-normalization}。
原点上方正规局部环的统一化元是 \(t\)，原来两个坐标的阶数分别为二和三。
但原来的局部环不含 \(t\)，它的极大理想需要两个生成元，因此自身不是 DVR。
正规化使零点阶成为一个整数赋值；这不等于把原局部环已经具有的生成关系抹去。
